\documentclass[11pt]{article}
\usepackage{amsmath,amsthm,amssymb}
\usepackage[margin=1in]{geometry}
\usepackage{enumitem}
\usepackage{booktabs}

\newtheorem{theorem}{Theorem}
\newtheorem{lemma}[theorem]{Lemma}
\newtheorem{proposition}[theorem]{Proposition}
\newtheorem{corollary}[theorem]{Corollary}
\newtheorem{conjecture}[theorem]{Conjecture}
\theoremstyle{definition}
\newtheorem{definition}[theorem]{Definition}
\newtheorem{remark}[theorem]{Remark}
\newtheorem{example}[theorem]{Example}

\title{The $j$-formula for $\lambda=(2,2,2,1^{n-6})$\\
  via branching reduction to $(2,2,1^{n-6})$}
\author{Clio}
\date{15 May 2026}

\begin{document}
\maketitle

\begin{abstract}
We give a fully rigorous proof of the $j$-formula at $j=n-2$ for
$\lambda=(2,2,2,1^{n-6})$ at every $n\ge 9$, completing the May-14
sketch.  The strategy is a \emph{branching reduction}: the
$H_q(S_{n-2})$-equivariant isomorphism
$\Phi: V_\mu|_{n-2}\xrightarrow{\sim} E_{n-1}^+|_{V_\lambda|_n}$
(where $\mu=(2,2,1^{n-6})$) reduces three of the four staircase
factors of $B_{n-2}^{(\lambda)}$ to the May-13 $j$-formula at $\mu$,
leaving only two extra factors $(T_{n-3}{+}1)(T_{n-2}{+}1)$ to
analyze explicitly.  We obtain
\[
   B_{n-2}^{(\lambda)} V_\lambda \;=\;
   \mathbb{Q}(q)\cdot(T_{n-3}{+}1)(T_{n-2}{+}1)\,\Phi(u^{(\mu)}),
\]
where $u^{(\mu)}\in V_\mu$ is the May-13 generator of
$B_{n-3}^{(\mu)} V_\mu$.  A direct case analysis on the
eight supports of $\Phi(u^{(\mu)})$ shows the right-hand side is
1-dimensional, contained in $E_1^-$, and supported on the eight
SYTs of $S_3^{(n)}$ from the recursive set defined in
\texttt{2026-05-14-2a-1n-family.tex}.

As a corollary, $\Pi^{S_n}|_{V_{(2,2,2,1^{n-6})}}=0$ unconditionally
for every $n\ge 9$ --- the second non-hook family in the rank-zero
program, now fully structurally proved.
\end{abstract}

\section{Setup}\label{sec:setup}

We use the conventions of \texttt{2026-05-13-non-hook-22-1n.tex}
and \texttt{2026-05-14-2a-1n-family.tex} (which we abbreviate
\textbf{May-13} and \textbf{May-14}).  $H_q(S_n)$ is the Hecke
algebra over $\mathbb{Q}(q)$ with generators $T_1,\ldots,T_{n-1}$
satisfying $T_i^2=(q-1)T_i+q$ and the braid relations.  $V_\lambda$
is the seminormal cell module of $H_q(S_n)$ for $\lambda\vdash n$,
with basis $\{v_T\}$ indexed by standard Young tableaux (SYTs).
We write
\[
   R'_k=(T_{k-1}{+}1)(T_{k-2}{+}1)\cdots(T_1{+}1),
   \qquad
   B_j=R'_jR'_{j+1}\cdots R'_n,
   \qquad
   \Pi^{S_n}=B_2.
\]
$E_i^\pm$ is the $\pm$-eigenspace of $T_i$ (eigenvalues $q$ for $+$,
$-1$ for $-$).

Throughout, $n\ge 9$ and
\[
   \lambda:=(2,2,2,1^{n-6})\vdash n,
   \qquad
   \mu:=(2,2,1^{n-6})\vdash n-2.
\]
$\ell(\lambda)=n-3$, so the $j$-formula is at
$j=\ell+1=n-2$ and $B_{n-2}^{(\lambda)}=R'_{n-2}R'_{n-1}R'_n$.

\subsection{Seminormal basis of $V_\lambda$}

By May-14 Lemma~1.2, SYTs of $\lambda$ are in bijection with triples
$(a_1,a_2,a_3)$ with $2\le a_1<a_2<a_3\le n$ satisfying the validity
conditions $a_2>b_1$, $a_3>b_2$, where $b_i$ is the $i$-th smallest
element of $\{2,\ldots,n\}\setminus\{a_1,a_2,a_3\}$.  We write
$v_{(a_1,a_2,a_3)}:=v_T$ for the corresponding SYT.

By May-14 Lemma~1.5, $T_1$ acts diagonally on $V_\lambda$:
\[
   E_1^+=\mathrm{span}\{v_{(2,a_2,a_3)}:\text{valid}\},
   \quad
   E_1^-=\mathrm{span}\{v_{(a_1,a_2,a_3)}:a_1\ge 3,\text{ valid}\}.
\]

\subsection{Seminormal basis of $V_\mu$}

By May-13 Lemma~2.1, SYTs of $\mu=(2,2,1^{m-4})$ at $m$ letters are
in bijection with pairs $(\alpha,\gamma)$ satisfying $2\le\alpha<\gamma\le m$
and $(\alpha,\gamma)\ne(2,3)$.  We use the convention
$v_{(\alpha,\gamma)}^{(\mu)}:=v_{T_{(\alpha,\gamma)}}$ in
$V_\mu|_m$.  Here $m=n-2$, so $(\alpha,\gamma)$ ranges in
$\{2,\ldots,n-2\}^2$.

\subsection{Phase A and the May-13 $j$-formula}

We rely on two prior results:

\begin{theorem}[Phase A for $\lambda$, May-14 Proposition~3.1]\label{thm:phaseA}
$R'_n V_\lambda=E_{n-1}^+$, of dimension $(n-2)(n-5)/2$.
\end{theorem}

\begin{theorem}[May-13 main theorem; \texttt{2026-05-13-non-hook-22-1n.tex}
Theorem~2.5]\label{thm:may13}
For $m\ge 6$, $\mu=(2,2,1^{m-4})\vdash m$ at length $\ell(\mu)=m-2$
and $j(\mu)=m-1$,
\[
   B_{m-1}^{(\mu)} V_\mu
   = R'^{(\mu)}_{m-1}R'^{(\mu)}_m V_\mu
   \;\subseteq\; E_1^-|_{V_\mu},
\]
of dimension $1$, with support exactly on the $4$ SYTs
$\{(m{-}3,m{-}2),(m{-}3,m{-}1),(m{-}2,m),(m{-}1,m)\}$.
Moreover, by May-13 Phase A endpoint (Proposition~3.1 of that
paper), $R'^{(\mu)}_m V_\mu=E_{m-1}^+|_{V_\mu}$, of dimension $m-3$.
\end{theorem}

We apply Theorem~\ref{thm:may13} at $m=n-2$.  Setting
$\hat u^{(\mu)}\in V_\mu$ to be a fixed generator of
$B_{n-3}^{(\mu)} V_\mu$ (unique up to nonzero scalar), May-13
gives the explicit form
\begin{equation}\label{eq:u-mu-form}
   \hat u^{(\mu)}
   = A\,v_{(n-5,n-4)}^{(\mu)} + B\,v_{(n-5,n-3)}^{(\mu)}
   + C\,v_{(n-4,n-2)}^{(\mu)} + D\,v_{(n-3,n-2)}^{(\mu)},
   \qquad A,B,C,D\in\mathbb{Q}(q)^\times.
\end{equation}

\section{The main theorem}

\begin{theorem}[$j$-formula at $\lambda=(2,2,2,1^{n-6})$,
$n\ge 9$]\label{thm:main}
$B_{n-2}^{(\lambda)} V_\lambda=R'_{n-2}R'_{n-1}R'_n V_\lambda$ is
$1$-dimensional, contained in $E_1^-|_{V_\lambda}$, with support
exactly the $8$ SYTs of $S_3^{(n)}$:
\[
   S_3^{(n)}=\Big\{
   (n{-}5,n{-}4,n{-}3),\ (n{-}5,n{-}4,n{-}2),\
   (n{-}5,n{-}3,n{-}1),\ (n{-}5,n{-}2,n{-}1),
\]
\[
   \phantom{S_3^{(n)}=\Big\{}\
   (n{-}4,n{-}3,n),\ (n{-}4,n{-}2,n),\
   (n{-}3,n{-}1,n),\ (n{-}2,n{-}1,n)\Big\}.
\]
\end{theorem}

\begin{corollary}[Rank-zero, structurally]\label{cor:rank-zero}
$\Pi^{S_n}|_{V_{(2,2,2,1^{n-6})}}=0$ for every $n\ge 9$.
\end{corollary}

\begin{proof}
Theorem~\ref{thm:main} gives $B_{n-2}^{(\lambda)} V_\lambda\subseteq E_1^-$.
Factor $\Pi^{S_n}=R'_2\cdots R'_{n-3}\cdot B_{n-2}^{(\lambda)}$.
The leftmost factor $R'_{n-3}$ ends in $(T_1{+}1)$, which annihilates
$E_1^-$.
\end{proof}

The remainder of the paper proves Theorem~\ref{thm:main}.

\section{The branching isomorphism $\Phi$}\label{sec:Phi}

The subalgebra of $H_q(S_n)$ generated by $\{T_1,\ldots,T_{n-3}\}$
is canonically isomorphic to $H_q(S_{n-2})$ (permuting letters
$\{1,\ldots,n-2\}$).  Since $T_{n-1}$ commutes with each
$T_i$ for $i\le n-3$ (index gap $\ge 2$), the eigenspace $E_{n-1}^+$
is preserved by the $H_q(S_{n-2})$-action and is therefore an
$H_q(S_{n-2})$-module.

\begin{definition}[The map $\Phi$]\label{def:Phi}
For each SYT $T'$ of $\mu$ on letters $\{1,\ldots,n-2\}$ with
column-$2$ entries $(\alpha,\gamma)$, define two SYTs $T_A, T_B$
of $\lambda$ on letters $\{1,\ldots,n\}$:
\begin{itemize}[leftmargin=2em,topsep=0.3em,itemsep=0.2em]
\item $T_A$: extend $T'$ by placing $n-1$ at $(3,2)$ and $n$ at the
new bottom row $(n-3,1)$.
\item $T_B$: extend $T'$ by placing $n$ at $(3,2)$ and $n-1$ at
$(n-3,1)$.
\end{itemize}
Both are valid SYTs of $\lambda$ (since $n-1,n$ exceed all other
entries, the column- and row-increase conditions are trivially
satisfied).  Letting $d=c(n)-c(n-1)$ be the axial distance in $T_A$
(where $c(r,c)=c-r$ denotes content), we have
$d=(-1)-(4-n)=n-5\ge 4$ for $n\ge 9$.

Define
\[
   \Phi(v_{T'}^{(\mu)})\;:=\;
   \text{the $T_{n-1}$ $+q$-eigenvector in the $2$-block }
   \{v_{T_A},v_{T_B}\}\subset V_\lambda,
\]
normalized by setting the coefficient of $v_{T_A}$ to $1$ (i.e.,
$\Phi(v_{T'}^{(\mu)})=v_{T_A}+\tau(d)v_{T_B}$ with $\tau(d)$ the
Hoefsmit ratio at axial distance $d$, $\tau(d)\in\mathbb{Q}(q)^\times$
for $|d|\ge 2$).  Extend $\Phi$ linearly to all of $V_\mu|_{n-2}$.

We write $w_\Phi(\alpha,\gamma):=\Phi(v_{(\alpha,\gamma)}^{(\mu)})$.
\end{definition}

\begin{lemma}[$\Phi$ is $H_q(S_{n-2})$-equivariant]\label{lem:Phi-equiv}
For every $1\le i\le n-3$ and every $v\in V_\mu$,
$T_i\,\Phi(v)=\Phi(T_i\,v)$.
\end{lemma}

\begin{proof}
By linearity, suffices to check $v=v_{T'}^{(\mu)}$.  The positions of
$i,i{+}1$ in $T_A$ and $T_B$ are identical to the positions in $T'$
(since $i,i{+}1\le n-2$ and the cells of $T_A,T_B$ holding entries
$\le n-2$ agree with $T'$).  Three cases for the $T_i$-action on
$T'$:

\paragraph{(a) Same row.}  $T_iv_{T'}^{(\mu)}=qv_{T'}^{(\mu)}$.  Then
$T_iv_{T_A}=qv_{T_A}$ and $T_iv_{T_B}=qv_{T_B}$ (same positions of
$i,i{+}1$).  Hence $T_i\Phi(v_{T'}^{(\mu)})
=T_i(v_{T_A}+\tau v_{T_B})=q(v_{T_A}+\tau v_{T_B})=q\Phi(v_{T'}^{(\mu)})
=\Phi(T_iv_{T'}^{(\mu)})$.

\paragraph{(b) Same column.}  Symmetric to (a), eigenvalue $-1$.

\paragraph{(c) $T_i$ $2$-block on $V_\mu$ at axial distance $d_i$.}
The block partner is $s_iT'$.  The $H_q(S_n)$-action on $V_\lambda$
has $T_i$-$2$-blocks at the SYTs $\{T_A,s_iT_A\}$ and
$\{T_B,s_iT_B\}$, each at the \emph{same} axial distance $d_i$
(the cells of $i,i+1$ are at identical positions across the four
SYTs).  Since $T_i$ and $T_{n-1}$ commute, the four-dimensional
subspace spanned by $\{v_{T_A},v_{T_B},v_{s_iT_A},v_{s_iT_B}\}$ is
preserved by both, and the actions decouple:
\begin{itemize}[leftmargin=2em,topsep=0.3em,itemsep=0.2em]
\item $T_{n-1}$: $2$-blocks $\{v_{T_A},v_{T_B}\}$ and
$\{v_{s_iT_A},v_{s_iT_B}\}$, axial distance $n-5$ in each.
\item $T_i$: $2$-blocks $\{v_{T_A},v_{s_iT_A}\}$ and
$\{v_{T_B},v_{s_iT_B}\}$, axial distance $d_i$ in each.
\end{itemize}
By Hoefsmit, both $T_{n-1}$-$2$-blocks have the same coefficient
matrix (same axial distance), so $\Phi(v_{T'}^{(\mu)})
=v_{T_A}+\tau v_{T_B}$ and
$\Phi(v_{s_iT'}^{(\mu)})=v_{s_iT_A}+\tau v_{s_iT_B}$ with the
same $\tau$.  Similarly both $T_i$-$2$-blocks have the same
coefficient matrix.  Write
$T_iv_{T_A}=a v_{T_A}+b v_{s_iT_A}$ and
$T_iv_{T_B}=a v_{T_B}+b v_{s_iT_B}$ (same $a,b$).  Then
\begin{align*}
   T_i\Phi(v_{T'}^{(\mu)})
   &= T_i(v_{T_A}+\tau v_{T_B})
   = a v_{T_A}+b v_{s_iT_A}+\tau(av_{T_B}+bv_{s_iT_B}) \\
   &= a(v_{T_A}+\tau v_{T_B})+b(v_{s_iT_A}+\tau v_{s_iT_B}) \\
   &= a\,\Phi(v_{T'}^{(\mu)})+b\,\Phi(v_{s_iT'}^{(\mu)})
   = \Phi(av_{T'}^{(\mu)}+bv_{s_iT'}^{(\mu)})
   = \Phi(T_iv_{T'}^{(\mu)}). \qedhere
\end{align*}
\end{proof}

\begin{lemma}[$\Phi$ is an iso onto $E_{n-1}^+$]\label{lem:Phi-iso}
$\Phi: V_\mu|_{n-2}\xrightarrow{\sim} E_{n-1}^+|_{V_\lambda|_n}$ is a
$\mathbb{Q}(q)$-linear isomorphism.
\end{lemma}

\begin{proof}
\emph{Injectivity.}  Distinct $T',T''$ of $\mu$ give disjoint pairs
$\{T_A,T_B\}$ and $\{T'_A,T'_B\}$ (the $\mu$-shape is recovered by
deleting the cells of $n-1,n$, which is unambiguous).  Hence
$\Phi(v_{T'}^{(\mu)})$ and $\Phi(v_{T''}^{(\mu)})$ have disjoint
seminormal supports, and the family $\{\Phi(v_{T'}^{(\mu)})\}_{T'}$
is linearly independent.

\emph{Image contained in $E_{n-1}^+$.}  Each $\Phi(v_{T'}^{(\mu)})$
is by construction a $T_{n-1}$ $+q$-eigenvector.

\emph{Dimension count.}  $\dim V_\mu=\frac{m(m-3)}{2}=\frac{(n-2)(n-5)}{2}$
(May-13 Lemma~2.4), equal to $\dim E_{n-1}^+|_{V_\lambda}$ by
Theorem~\ref{thm:phaseA}.
\end{proof}

\begin{lemma}[$\Phi$ preserves $E_1^-$]\label{lem:Phi-E1}
$\Phi(E_1^-|_{V_\mu})\subseteq E_1^-|_{V_\lambda}$, and
$\Phi(E_1^+|_{V_\mu})\subseteq E_1^+|_{V_\lambda}$.
\end{lemma}

\begin{proof}
$T_1$ is in the $H_q(S_{n-2})$-subalgebra and $\Phi$ is
$T_1$-equivariant (Lemma~\ref{lem:Phi-equiv}).  Hence
$T_1v=\pm v$ implies $T_1\Phi(v)=\pm\Phi(v)$.
\end{proof}

\section{The reduction formula}\label{sec:reduction}

\begin{theorem}[Reduction]\label{thm:reduction}
Set $S:=(T_{n-4}{+}1)(T_{n-5}{+}1)\cdots(T_1{+}1)\in H_q(S_{n-3})$,
and recall $R'_{n-2}=(T_{n-3}{+}1)\,S$.  Then
\begin{equation}\label{eq:reduction}
   B_{n-2}^{(\lambda)}\,V_\lambda
   \;=\;
   (T_{n-3}{+}1)(T_{n-2}{+}1)\,\Phi\big(B_{n-3}^{(\mu)} V_\mu\big).
\end{equation}
In particular, $\dim B_{n-2}^{(\lambda)}V_\lambda\le 1$.
\end{theorem}

\begin{proof}
We compute step by step.

\emph{Step 1.}  $R'_n V_\lambda=E_{n-1}^+=\Phi(V_\mu)$ by
Theorem~\ref{thm:phaseA} and Lemma~\ref{lem:Phi-iso}.

\emph{Step 2.}  $R'_{n-2}\in H_q(S_{n-2})$ uses only $T_1,\ldots,T_{n-3}$.
By Lemma~\ref{lem:Phi-equiv}, $\Phi$ intertwines this action:
\[
   R'_{n-2}\,\Phi(V_\mu)=\Phi(R'^{(\mu)}_{n-2} V_\mu)
   =\Phi(E_{n-3}^+|_{V_\mu}),
\]
the last equality by May-13 Phase A (Theorem~\ref{thm:may13}).

\emph{Step 3.}  $R'_{n-1}=(T_{n-2}{+}1)R'_{n-2}$, so
\[
   R'_{n-1}R'_n V_\lambda=(T_{n-2}{+}1)\,R'_{n-2}\,E_{n-1}^+
   =(T_{n-2}{+}1)\,\Phi(E_{n-3}^+|_{V_\mu}).
\]

\emph{Step 4.}  $R'_{n-2}=(T_{n-3}{+}1)S$ as in the theorem statement.
Since $S$ uses generators $T_1,\ldots,T_{n-4}$ and $T_{n-2}$ has
index gap $\ge 2$ from each, $S$ commutes with $T_{n-2}$, hence with
$(T_{n-2}{+}1)$.  Therefore
\[
   R'_{n-2}\,(T_{n-2}{+}1)
   = (T_{n-3}{+}1)\,S\,(T_{n-2}{+}1)
   = (T_{n-3}{+}1)\,(T_{n-2}{+}1)\,S.
\]

\emph{Step 5.}  $S\in H_q(S_{n-3})\subset H_q(S_{n-2})$ is
$\Phi$-equivariant (Lemma~\ref{lem:Phi-equiv}), and
$S=R'^{(\mu)}_{n-3}$ as an element of $H_q(S_{n-2})$.  Hence
\[
   S\,\Phi(E_{n-3}^+|_{V_\mu})
   = \Phi(R'^{(\mu)}_{n-3}\,E_{n-3}^+|_{V_\mu})
   = \Phi(R'^{(\mu)}_{n-3}R'^{(\mu)}_{n-2} V_\mu)
   = \Phi(B_{n-3}^{(\mu)} V_\mu).
\]

\emph{Combining.}  Therefore
\begin{align*}
   B_{n-2}^{(\lambda)} V_\lambda
   &= R'_{n-2}\,R'_{n-1}\,R'_n V_\lambda
   = R'_{n-2}\,(T_{n-2}{+}1)\,\Phi(E_{n-3}^+|_{V_\mu}) \\
   &= (T_{n-3}{+}1)(T_{n-2}{+}1)\,S\,\Phi(E_{n-3}^+|_{V_\mu}) \\
   &= (T_{n-3}{+}1)(T_{n-2}{+}1)\,\Phi(B_{n-3}^{(\mu)} V_\mu).
\end{align*}
By Theorem~\ref{thm:may13}, $\dim B_{n-3}^{(\mu)} V_\mu=1$, so
$\dim\Phi(B_{n-3}^{(\mu)}V_\mu)=1$ (as $\Phi$ is injective) and
the right-hand side has dimension at most $1$.
\end{proof}

\begin{corollary}[$E_1^-$ containment]\label{cor:E1-containment}
$B_{n-2}^{(\lambda)} V_\lambda\subseteq E_1^-|_{V_\lambda}$.
\end{corollary}

\begin{proof}
By May-13, $B_{n-3}^{(\mu)}V_\mu\subseteq E_1^-|_{V_\mu}$.  By
Lemma~\ref{lem:Phi-E1}, $\Phi(B_{n-3}^{(\mu)}V_\mu)\subseteq
E_1^-|_{V_\lambda}$.  The generators $T_{n-3}$ and $T_{n-2}$
commute with $T_1$ (index gap $\ge 2$ for $n\ge 6$), so
$(T_{n-3}{+}1)(T_{n-2}{+}1)$ preserves each $T_1$-eigenspace.
In particular, $E_1^-$ is preserved.  Apply
Theorem~\ref{thm:reduction}.
\end{proof}

\section{Computing the lifted image}\label{sec:lift}

To complete the proof of Theorem~\ref{thm:main}, we must show:
\begin{enumerate}[leftmargin=2em,topsep=0.3em,itemsep=0.2em]
\item $(T_{n-3}{+}1)(T_{n-2}{+}1)\,\Phi(\hat u^{(\mu)})\ne 0$
(so the image is exactly $1$-dimensional, not $0$).
\item Its seminormal support equals $S_3^{(n)}$ from
Theorem~\ref{thm:main}, with all $8$ coefficients nonzero.
\end{enumerate}

Recall \eqref{eq:u-mu-form}: $\hat u^{(\mu)}=
A v_{(n-5,n-4)}^{(\mu)}+B v_{(n-5,n-3)}^{(\mu)}
+C v_{(n-4,n-2)}^{(\mu)}+D v_{(n-3,n-2)}^{(\mu)}$, with
$A,B,C,D\ne 0$.  Applying $\Phi$:
\begin{equation}\label{eq:Phi-u-mu}
   \Phi(\hat u^{(\mu)})
   = A\,w_\Phi(n{-}5,n{-}4)+B\,w_\Phi(n{-}5,n{-}3)
   + C\,w_\Phi(n{-}4,n{-}2)+D\,w_\Phi(n{-}3,n{-}2),
\end{equation}
where each $w_\Phi(\alpha,\gamma)=v_{(\alpha,\gamma,n-1)}
+\tau v_{(\alpha,\gamma,n)}$ ($\tau\in\mathbb{Q}(q)^\times$ at axial
distance $n-5\ge 4$).

\subsection{Step (a): apply $(T_{n-2}{+}1)$}

We compute $T_{n-2}$ on each of the $8$ seminormal supports of
$\Phi(\hat u^{(\mu)})$, namely $v_{(\alpha,\gamma,\delta)}$ for
$(\alpha,\gamma)\in S_2^{(n-2)}$ and $\delta\in\{n-1,n\}$.  Note
$T_{n-2}$ involves the entries $n-2$ and $n-1$.

\begin{lemma}[Case analysis of $T_{n-2}$]\label{lem:Tn2-cases}
For $(\alpha,\gamma)\in S_2^{(n-2)}$ and $\delta\in\{n-1,n\}$:
\begin{enumerate}[leftmargin=2em,topsep=0.3em,itemsep=0.2em,
label=\textup{(\alph*)}]
\item $(\alpha,\gamma,\delta)=(n{-}5,n{-}4,n)$ or
$(n{-}5,n{-}3,n)$: $n-2$ and $n-1$ are both in column~$1$ at
adjacent rows.  $T_{n-2}v=-v$.
\item $(\alpha,\gamma,\delta)=(n{-}4,n{-}2,n{-}1)$ or
$(n{-}3,n{-}2,n{-}1)$: $n-2,n-1$ are at $(2,2),(3,2)$ (same
column~$2$).  $T_{n-2}v=-v$.
\item $(\alpha,\gamma,\delta)=(n{-}5,n{-}4,n{-}1)$ or
$(n{-}5,n{-}3,n{-}1)$: $n-1$ at $(3,2)$, $n-2$ at $(n{-}4,1)$.
$T_{n-2}$ acts as a $2$-block at axial distance $d=n-6$ coupling
$v_{(\alpha,\gamma,n-1)}\leftrightarrow v_{(\alpha,\gamma,n-2)}$.
\item $(\alpha,\gamma,\delta)=(n{-}4,n{-}2,n)$ or
$(n{-}3,n{-}2,n)$: $n-2$ at $(2,2)$ (content $0$), $n-1$ at
$(n{-}3,1)$ (content $4-n$).  $T_{n-2}$ acts as a $2$-block at
axial distance $d=4-n$ coupling
$v_{(\alpha,n-2,n)}\leftrightarrow v_{(\alpha,n-1,n)}$ (where the
second column-$2$ entry shifts from $n-2$ to $n-1$).
\end{enumerate}
For $n\ge 9$, all axial distances in (c), (d) satisfy
$|d|\ge 3$, hence the $2$-blocks are non-degenerate
(Hoefsmit coefficients in $\mathbb{Q}(q)^\times$).
\end{lemma}

\begin{proof}
Direct case-by-case verification using Lemma~1.4 of May-14
(positions in seminormal SYTs).  For (a) with
$(\alpha,\gamma)=(n{-}5,n{-}4)$, $\delta=n$:  column-$1$ entries
are $\{1\}\cup\{2,\ldots,n\}\setminus\{n{-}5,n{-}4,n\}$, sorted
$1,2,\ldots,n{-}6,n{-}3,n{-}2,n{-}1$; thus $n-2$ at row $n-4$,
$n-1$ at row $n-3$, adjacent.  Other cases similar.

For the swap-partner SYTs in (c), (d): e.g., in (c) with
$(\alpha,\gamma)=(n{-}5,n{-}4)$, swapping $n{-}2$ and $n{-}1$
moves $n{-}1$ to row $n-4$ of column~$1$ and $n{-}2$ to $(3,2)$,
yielding SYT $(n{-}5,n{-}4,n{-}2)$, which is valid since
$a_3=n{-}2>b_2=3$ for $n\ge 6$.

The axial distances follow from $c(i)=\mathrm{col}(i)-\mathrm{row}(i)$
applied to each case: e.g., $(c(n{-}1)-c(n{-}2))$ in (c) is
$(-1)-((1-(n-4)))=(-1)-(5-n)=n-6$.

For $n\ge 9$: $|d|\ge|n-6|\ge 3$ and $|4-n|\ge 5$, both $\ne 1$.
\end{proof}

\begin{proposition}[Image after $(T_{n-2}{+}1)$]\label{prop:Tn2-image}
$(T_{n-2}{+}1)\,\Phi(\hat u^{(\mu)})$ is a linear combination of
the following $8$ seminormal basis vectors, with all coefficients
nonzero:
\begin{align*}
\text{from }w_\Phi(n{-}5,n{-}4):\quad &
v_{(n-5,n-4,n-1)},\ v_{(n-5,n-4,n-2)}\\
\text{from }w_\Phi(n{-}5,n{-}3):\quad &
v_{(n-5,n-3,n-1)},\ v_{(n-5,n-3,n-2)}\\
\text{from }w_\Phi(n{-}4,n{-}2):\quad &
v_{(n-4,n-2,n)},\ v_{(n-4,n-1,n)}\\
\text{from }w_\Phi(n{-}3,n{-}2):\quad &
v_{(n-3,n-2,n)},\ v_{(n-3,n-1,n)}
\end{align*}
The eight supports are pairwise disjoint.
\end{proposition}

\begin{proof}
For each of the $4$ terms in \eqref{eq:Phi-u-mu}, apply
$(T_{n-2}{+}1)$ to the two-element seminormal expansion of
$w_\Phi(\alpha,\gamma)=v_{(\alpha,\gamma,n-1)}+\tau
v_{(\alpha,\gamma,n)}$.  By Lemma~\ref{lem:Tn2-cases}:

\emph{Term 1: $w_\Phi(n{-}5,n{-}4)$.}  By Lemma~\ref{lem:Tn2-cases}(a),
$(T_{n-2}{+}1)v_{(n-5,n-4,n)}=0$.  By Lemma~\ref{lem:Tn2-cases}(c),
$(T_{n-2}{+}1)v_{(n-5,n-4,n-1)}=\kappa\big(\alpha' v_{(n-5,n-4,n-1)}
+\beta' v_{(n-5,n-4,n-2)}\big)$ with
$\kappa,\alpha',\beta'\in\mathbb{Q}(q)^\times$ (Hoefsmit at $d=n-6$).
The resulting support is $\{(n-5,n-4,n-1),(n-5,n-4,n-2)\}$, with
the coefficient $A\kappa\alpha'\ne 0$ at the first SYT and
$A\kappa\beta'\ne 0$ at the second.

\emph{Term 2: $w_\Phi(n{-}5,n{-}3)$.}  By Lemma~\ref{lem:Tn2-cases}(a),
$(T_{n-2}{+}1)v_{(n-5,n-3,n)}=0$.  By Lemma~\ref{lem:Tn2-cases}(c),
support of $(T_{n-2}{+}1)v_{(n-5,n-3,n-1)}$ is
$\{(n-5,n-3,n-1),(n-5,n-3,n-2)\}$.  Coefficients $B\kappa\alpha'$
and $B\kappa\beta'$, both nonzero.

\emph{Term 3: $w_\Phi(n{-}4,n{-}2)$.}  By Lemma~\ref{lem:Tn2-cases}(b),
$(T_{n-2}{+}1)v_{(n-4,n-2,n-1)}=0$.  By Lemma~\ref{lem:Tn2-cases}(d),
$(T_{n-2}{+}1)v_{(n-4,n-2,n)}=\kappa'\big(\alpha'' v_{(n-4,n-2,n)}
+\beta'' v_{(n-4,n-1,n)}\big)$ with
$\kappa',\alpha'',\beta''\in\mathbb{Q}(q)^\times$ (Hoefsmit at
$d=4-n$).  Support $\{(n-4,n-2,n),(n-4,n-1,n)\}$.  Coefficients
$C\tau\kappa'\alpha''$ and $C\tau\kappa'\beta''$ (the $\tau$
from $w_\Phi$), both nonzero.

\emph{Term 4: $w_\Phi(n{-}3,n{-}2)$.}  Same as Term 3 with
$\alpha=n-3$.  Support $\{(n-3,n-2,n),(n-3,n-1,n)\}$.

\emph{Disjointness.}  The $8$ supports listed all have distinct
$(\alpha,\gamma,\delta)$-triples (different middle column-$2$
entries or different positions of the boundary index), so they
are pairwise disjoint.  Hence the $8$ coefficients are
independent and all nonzero.
\end{proof}

\subsection{Step (b): apply $(T_{n-3}{+}1)$}

\begin{lemma}[Case analysis of $T_{n-3}$]\label{lem:Tn3-cases}
For each of the $8$ SYTs in Proposition~\ref{prop:Tn2-image}:
\begin{enumerate}[leftmargin=2em,topsep=0.3em,itemsep=0.2em,
label=\textup{(\alph*)}]
\item $(n{-}5,n{-}4,n{-}1)$: $n-3$ at $(n-5,1)$, $n-2$ at $(n-4,1)$.
Same column~$1$ adjacent.  $T_{n-3}v=-v$.
\item $(n{-}5,n{-}4,n{-}2)$: $n-3$ at $(n-5,1)$ (content $6-n$),
$n-2$ at $(3,2)$ (content $-1$).  $T_{n-3}$ acts as a $2$-block at
axial distance $d=n-7$, coupling
$v_{(n-5,n-4,n-2)}\leftrightarrow v_{(n-5,n-4,n-3)}$.
\item $(n{-}5,n{-}3,n{-}1)$: $n-3$ at $(2,2)$ (content $0$),
$n-2$ at $(n-4,1)$ (content $5-n$).  $T_{n-3}$ acts as a $2$-block
at axial distance $d=5-n$, coupling
$v_{(n-5,n-3,n-1)}\leftrightarrow v_{(n-5,n-2,n-1)}$.
\item $(n{-}5,n{-}3,n{-}2)$: $n-3,n-2$ at $(2,2),(3,2)$ (same
column~$2$ adjacent).  $T_{n-3}v=-v$.
\item $(n{-}4,n{-}2,n)$: $n-3$ at $(n-4,1)$ (content $5-n$), $n-2$
at $(2,2)$ (content $0$).  $T_{n-3}$ acts as a $2$-block at axial
distance $d=n-5$, coupling
$v_{(n-4,n-2,n)}\leftrightarrow v_{(n-4,n-3,n)}$.
\item $(n{-}4,n{-}1,n)$: $n-3$ at $(n-4,1)$, $n-2$ at $(n-3,1)$.
Same column~$1$ adjacent.  $T_{n-3}v=-v$.
\item $(n{-}3,n{-}2,n)$: $n-3,n-2$ at $(1,2),(2,2)$ (same
column~$2$ adjacent).  $T_{n-3}v=-v$.
\item $(n{-}3,n{-}1,n)$: $n-3$ at $(1,2)$ (content $1$), $n-2$ at
$(n-3,1)$ (content $4-n$).  $T_{n-3}$ acts as a $2$-block at axial
distance $d=3-n$, coupling
$v_{(n-3,n-1,n)}\leftrightarrow v_{(n-2,n-1,n)}$.
\end{enumerate}
For $n\ge 9$, all axial distances in (b), (c), (e), (h) satisfy
$|d|\ge 2$, so $2$-blocks are non-degenerate.
\end{lemma}

\begin{proof}
Direct case analysis using the position formulas for the
seminormal basis of $\lambda$.  Cases (a), (d), (f), (g) are
``$T_{n-3}=-1$'' cases (same-column adjacency); cases (b), (c),
(e), (h) are $2$-blocks.  For $n=9$: $|n-7|=2$, $|5-n|=4$,
$|n-5|=4$, $|3-n|=6$, all $\ge 2$.

For the swap-partner SYTs: e.g., in (b), swap $n{-}3$ and $n{-}2$
in $(n{-}5,n{-}4,n{-}2)$ moves $n{-}3$ to $(3,2)$ and $n{-}2$ to
$(n{-}5,1)$, yielding $(n{-}5,n{-}4,n{-}3)$ — valid since
$n{-}4<n{-}3<a_2=n{-}4$ wait no: $a_2=n{-}4$ and the new $a_3=n{-}3$
satisfies $a_2<a_3$ since $n{-}4<n{-}3$.  Validity of column~$1$:
$b_2=3<n{-}3$.  Good.

In (c), swap moves $n{-}2$ to $(2,2)$ and $n{-}3$ to $(n{-}4,1)$,
yielding $(n{-}5,n{-}2,n{-}1)$ — valid since
$n{-}5<n{-}2<n{-}1$ and $b_1=2<a_2=n{-}2$.

In (e), swap moves $n{-}3$ to $(2,2)$ and $n{-}2$ to $(n{-}4,1)$,
yielding $(n{-}4,n{-}3,n)$ — valid since $n{-}4<n{-}3<n$ and
$b_1=2<n{-}3$.

In (h), swap moves $n{-}2$ to $(1,2)$ and $n{-}3$ to $(n{-}3,1)$,
yielding $(n{-}2,n{-}1,n)$ — valid since
$n{-}2<n{-}1<n$ and $b_1=2<n{-}2$.
\end{proof}

\begin{proposition}[Final support]\label{prop:final}
$(T_{n-3}{+}1)(T_{n-2}{+}1)\,\Phi(\hat u^{(\mu)})$ has seminormal
support exactly equal to the $8$-element set $S_3^{(n)}$ of
Theorem~\ref{thm:main}, with all $8$ coefficients nonzero.
\end{proposition}

\begin{proof}
By Proposition~\ref{prop:Tn2-image}, $(T_{n-2}{+}1)\Phi(\hat u^{(\mu)})$
has $8$-element pairwise-disjoint support:
\begin{align*}
   \mathrm{supp}_1&=\{(n{-}5,n{-}4,n{-}1),(n{-}5,n{-}4,n{-}2)\},
   &
   \mathrm{supp}_2&=\{(n{-}5,n{-}3,n{-}1),(n{-}5,n{-}3,n{-}2)\},\\
   \mathrm{supp}_3&=\{(n{-}4,n{-}2,n),(n{-}4,n{-}1,n)\},
   &
   \mathrm{supp}_4&=\{(n{-}3,n{-}2,n),(n{-}3,n{-}1,n)\}.
\end{align*}
Each $\mathrm{supp}_i$ has both coefficients nonzero.

Apply $(T_{n-3}{+}1)$:

\emph{$\mathrm{supp}_1$: $(n{-}5,n{-}4,n{-}1)$ killed
(Lemma~\ref{lem:Tn3-cases}(a)); $(n{-}5,n{-}4,n{-}2)$ produces
$2$-block to $(n{-}5,n{-}4,n{-}3)$ (Lemma~\ref{lem:Tn3-cases}(b)).}
New support: $\{(n{-}5,n{-}4,n{-}2),(n{-}5,n{-}4,n{-}3)\}$, two
nonzero coefficients.

\emph{$\mathrm{supp}_2$: $(n{-}5,n{-}3,n{-}1)$ produces $2$-block to
$(n{-}5,n{-}2,n{-}1)$ (Lemma~\ref{lem:Tn3-cases}(c));
$(n{-}5,n{-}3,n{-}2)$ killed (Lemma~\ref{lem:Tn3-cases}(d)).}
New support: $\{(n{-}5,n{-}3,n{-}1),(n{-}5,n{-}2,n{-}1)\}$.

\emph{$\mathrm{supp}_3$: $(n{-}4,n{-}2,n)$ produces $2$-block to
$(n{-}4,n{-}3,n)$ (Lemma~\ref{lem:Tn3-cases}(e));
$(n{-}4,n{-}1,n)$ killed (Lemma~\ref{lem:Tn3-cases}(f)).}
New support: $\{(n{-}4,n{-}2,n),(n{-}4,n{-}3,n)\}$.

\emph{$\mathrm{supp}_4$: $(n{-}3,n{-}2,n)$ killed
(Lemma~\ref{lem:Tn3-cases}(g)); $(n{-}3,n{-}1,n)$ produces $2$-block
to $(n{-}2,n{-}1,n)$ (Lemma~\ref{lem:Tn3-cases}(h)).}
New support: $\{(n{-}3,n{-}1,n),(n{-}2,n{-}1,n)\}$.

Combining: the eight final supports are
\[
   \{(n{-}5,n{-}4,n{-}3),(n{-}5,n{-}4,n{-}2),
   (n{-}5,n{-}3,n{-}1),(n{-}5,n{-}2,n{-}1),
\]
\[
   (n{-}4,n{-}3,n),(n{-}4,n{-}2,n),
   (n{-}3,n{-}1,n),(n{-}2,n{-}1,n)\},
\]
which is exactly $S_3^{(n)}$.

\emph{Nonzero coefficients.}  We track coefficients: each
$\mathrm{supp}_i$ contributes two new SYTs with coefficients
of the form
$(\text{nonzero coeff of }(T_{n-2}{+}1)\text{ step})\cdot\kappa''\cdot
(\alpha'''\text{ or }\beta''')$, where $\kappa'',\alpha''',\beta'''$
are the Hoefsmit coefficients at the $T_{n-3}$ axial distance, all
nonzero.  The $8$ resulting supports are pairwise disjoint (no two
sub-supports overlap, by inspection of the lists).  Hence the
$8$ coefficients are independent products of nonzero terms, all
nonzero.
\end{proof}

\subsection{Proof of Theorem~\ref{thm:main}}

\begin{proof}[Proof of Theorem~\ref{thm:main}]
By Theorem~\ref{thm:reduction},
$B_{n-2}^{(\lambda)}V_\lambda=(T_{n-3}{+}1)(T_{n-2}{+}1)\Phi(\hat u^{(\mu)})$
up to scalar.  By Proposition~\ref{prop:final}, this is nonzero
and has support exactly $S_3^{(n)}$ with $8$ nonzero coefficients.
Hence $\dim B_{n-2}^{(\lambda)}V_\lambda=1$.

The $E_1^-$ containment is Corollary~\ref{cor:E1-containment} (and
is also immediate since every SYT in $S_3^{(n)}$ has minimum entry
$\ge n-5\ge 4>2$, so $2$ is in column~$1$ and the basis vector lies
in $E_1^-$).
\end{proof}

\section{Computational verification}\label{sec:verify}

The proof above is rigorous for all $n\ge 9$.  We have additionally
verified it computationally in symbolic $\mathbb{Q}(q)$ at $n=9$
and at the specific rational $q=7$ for $n\in\{9,10,11\}$ (the
verification script is
\texttt{\textasciitilde/projects/scratch/2026-05-15-a3-reduction/verify\_reduction.py}
and \texttt{verify\_symbolic.py}).

At $n=9$, the explicit form of the generator of
$B_7^{(\lambda)}V_\lambda$ (in the seminormal basis, normalized
with first nonzero coefficient set to $1$) is, with all
coefficients nonzero rational functions in $q$:
\begin{align*}
   \mathrm{coef}\,v_{(4,5,6)} &= q^9(q+1)^2/(q^2+1), \\
   \mathrm{coef}\,v_{(4,5,7)} &= q^8(q+1)^3/(q^2+1), \\
   \mathrm{coef}\,v_{(4,6,8)} &= q^7(q+1)(q^4+2q^3+3q^2+2q+1)/(q^2+1)^2, \\
   \mathrm{coef}\,v_{(4,7,8)} &= q^6(q+1)(q^3+2q^2+2q+1)/(q^2+1), \\
   \mathrm{coef}\,v_{(5,6,9)} &= q^6(q+1)(q^4+3q^3+4q^2+3q+1)/(q^4+q^3+q^2+q+1), \\
   \mathrm{coef}\,v_{(5,7,9)} &= q^5\cdot\text{(nonzero in }\mathbb{Q}(q)), \\
   \mathrm{coef}\,v_{(6,8,9)} &= q^4\cdot\text{(nonzero in }\mathbb{Q}(q)), \\
   \mathrm{coef}\,v_{(7,8,9)} &= q^3(q+1)^3.
\end{align*}
The support is exactly the $8$ SYTs of $S_3^{(9)}$, as predicted.

The same $8$-element support pattern (shifted by $n-5$) appears at
$n=10,11$.

\section{Discussion}\label{sec:discuss}

\subsection{The structural picture}

Theorem~\ref{thm:reduction} expresses $B_{n-2}^{(\lambda)}V_\lambda$
as a single $H_q(S_{n-2})$-equivariant lift of $B_{n-3}^{(\mu)}V_\mu$
followed by $(T_{n-3}{+}1)(T_{n-2}{+}1)$.  This is the
representation-theoretic ``$T_2$-obstruction'' from the May-14 note
turned around: the obstruction to a clean reduction
$(a\text{ at }n)\mapsto(a-1\text{ at }n-2)$ via simply forgetting
the top row is now isolated to the two extra factors
$(T_{n-3}{+}1)(T_{n-2}{+}1)$, which can be analyzed case-by-case.

\subsection{Pattern for general $a$}

Conjecturally, the same strategy extends to $\lambda=(2^a,1^{n-2a})$
at $n\ge 3a$:
\begin{itemize}[leftmargin=2em,topsep=0.3em,itemsep=0.2em]
\item Phase A: $R'_n V_\lambda=E_{n-1}^+\cong V_{\mu_1}|_{n-2}$ via
$\Phi$, where $\mu_1=(2^{a-1},1^{n-2a})$.
\item Iterating Theorem~\ref{thm:reduction}-style commutations,
\[
   B_{n-a+1}^{(\lambda)}V_\lambda
   = \big[\prod_{k=0}^{a-2}(T_{n-a+1-k}{+}1)\big]\cdot
   \Phi^{(a)}\big(B_?^{(\mu_{a-1})}V_{\mu_{a-1}}\big),
\]
where $\mu_{a-1}$ is the partition reached by iterated branching
and $\Phi^{(a)}$ is the composite branching iso.
\item The product of $(T_j{+}1)$ factors performs the ``L/R
boundary choice'' at each level, generating the $2^a$ support via
the recursive structure $S_a$.
\end{itemize}
This is not yet proved.  The key obstruction for general $a$ is
verifying that each successive branching iso $\Phi^{(k)}$ exists
(needs the $(T_j{+}1)$-image after the $k$-th step to factor
through a single $H_q(S_{n-2k})$-isotypic component).  At $a=3$
the analogue holds because $\Phi(B_{n-3}^{(\mu)}V_\mu)\subseteq
E_{n-3}^+\cap E_{n-1}^+$ by Step~2 of Theorem~\ref{thm:reduction},
which is structurally a single isotypic copy.

\subsection{The $\mathbb{Z}/2$ symmetry of $S_a$}

The set $S_3$ (and generally $S_a$) is closed under the involution
$x\mapsto(2a-1)-x$ (May-14 Section~3).  This reversal, applied
to the $8$ SYTs of $S_3^{(n)}$, swaps:
\begin{align*}
   (n{-}5,n{-}4,n{-}3)&\leftrightarrow(n{-}2,n{-}1,n),\\
   (n{-}5,n{-}4,n{-}2)&\leftrightarrow(n{-}3,n{-}1,n),\\
   (n{-}5,n{-}3,n{-}1)&\leftrightarrow(n{-}4,n{-}2,n),\\
   (n{-}5,n{-}2,n{-}1)&\leftrightarrow(n{-}4,n{-}3,n).
\end{align*}
This is exactly the involution exchanging the ``$\mathrm{supp}_i$
side'' (post-$(T_{n-2}{+}1)$) of Proposition~\ref{prop:Tn2-image}:
$\mathrm{supp}_1\leftrightarrow\mathrm{supp}_4$,
$\mathrm{supp}_2\leftrightarrow\mathrm{supp}_3$ after one
$(T_{n-3}{+}1)$ application.

The representation-theoretic source: $\lambda=(2,2,2,1^{n-6})$ has
two ``extremal'' length-preserving corners (the column-$2$ corner
$(3,2)$ and the column-$1$ ``bottom'' which is conceptually
adjacent), and the L/R choice corresponds to which corner is
``filled first'' as one fills cells from $n$ downward.

\section{Honest assessment}

\paragraph{What is fully proved.}
\begin{enumerate}[leftmargin=2em,topsep=0.3em,itemsep=0.2em]
\item Theorem~\ref{thm:main}: the $j$-formula for
$\lambda=(2,2,2,1^{n-6})$ at every $n\ge 9$ is rigorously
established, with support exactly $S_3^{(n)}$ ($8$ SYTs, all
coefficients nonzero in $\mathbb{Q}(q)$).
\item Corollary~\ref{cor:rank-zero}: rank-zero
$\Pi^{S_n}|_{V_{(2,2,2,1^{n-6})}}=0$ for all $n\ge 9$, the second
non-hook family fully structurally proved.
\item The branching reduction Theorem~\ref{thm:reduction} is a
clean structural statement, isolating the entire $a=3$ proof to
the May-13 $a=2$ result plus two explicit Hoefsmit case analyses.
\end{enumerate}

\paragraph{What is left open.}
\begin{enumerate}[leftmargin=2em,topsep=0.3em,itemsep=0.2em]
\item The general-$a$ conjecture (May-14 Conjecture~5.1) remains
conjectural.  The reduction-via-branching strategy here suggests a
path: iterate Theorem~\ref{thm:reduction}-style commutations, but
verifying that each successive branching iso lands in a single
isotypic component is nontrivial.
\item The combinatorial origin of $S_a$ (Dyck-like paths? parking
functions? Eulerian numbers?) is not identified.  The $\mathbb{Z}/2$
reversal symmetry suggests palindromic combinatorics.
\end{enumerate}

\paragraph{Status of the rank-zero program.}
\begin{itemize}[leftmargin=2em,topsep=0.3em,itemsep=0.2em]
\item Fully structurally proved: $(1^n)$, $(2,1^{n-2})$,
$(3,1^{n-3})$ at $n\ge 5$, $(2,2,1^{n-4})$ at $n\ge 6$, and now
$(2,2,2,1^{n-6})$ at $n\ge 9$.
\item Computationally verified: $(2^a,1^{n-2a})$ at $a=4,5$ and
$n\ge 3a$, support $=2^a$.
\item Conjectural: general $\lambda=(2^a,1^{n-2a})$ at $n\ge 3a$.
\end{itemize}

\end{document}
